\chapter{Metodología de la Investigación}
\section{Tipo de Investigación}
    \begin{sectext}

        El tipo de investigación enmarca la idea del qué hacer y qué información necesito para lograr ese objetivo. Cada tipo de investigación conlleva una secuencia estructurada de fases y operaciones. Según \textcite{carrasco_dias_metodologia_2005}, existen 4 tipos de investigación: Investigación básica, donde se busca ampliar el conocimiento científico existente acerca de una realidad, el objetivo de este tipo de investigaciones recae sobre las teorías científicas; Investigación Aplicada, donde se buscan resultados prácticos e inmediatos por medio de actuar, transformar o modificar una parte de la realidad en estudio; Investigación Sustantiva, donde se busca responder a interrogantes concretas en favor de la consolidación de las teorías científicas y la aplicación de las investigaciones aplicadas y tecnológicas; por último, la Investigación Tecnológica busca perfeccionar u manipular un sector de la realidad en estudio mediante la aplicación de nuevos sistemas, modelos u técnicas.
        
        El presente trabajo se enmarca en el tipo de investigación tecnológica debido al objetivo central de realizar una invención tecnológica...
    \end{sectext}

\section{Nivel de la Investigación}
    \begin{sectext}
        El nivel de la investigación establece la profundidad que tendrá un estudio respecto a la problemática identificada. Según \textcite{carrasco_dias_metodologia_2005}, los niveles de investigación representan el orden secuencial y coherente que deben tener las investigaciones; se catalogan 4 tipos: preliminares o exploratorios, donde se establecen las primeras ideas y conceptos de la realidad a investigar; descriptivos, donde establecemos de manera robusta las propiedades, características, objetos, cualidades, rasgos u fenómenos de la realidad a estudiar; explicativos o causales, donde establecemos las causas u factores que dan origen a las circunstancias y naturaleza del fenómeno en estudio; y experimentales, en el cual se establece un nuevo paradigma, sistema, modelo u tratamiento para intentar corregir la situación problemática de la realidad en estudio.
        
        El presente trabajo se enmarca en el nivel de investigación de tipo experimental debido a que se busca entender cómo un nuevo sistema...
    \end{sectext}